\documentclass[11pt]{article}
\usepackage[margin=1in]{geometry}
\usepackage{amsmath,amssymb}
\usepackage{hyperref}
\usepackage{times}

\hypersetup{colorlinks=true,linkcolor=blue,citecolor=blue,urlcolor=blue}

\begin{document}

\noindent
\textbf{To:} \\
Editor-in-Chief \\
\textit{Computers \& Security} \\
Elsevier

\vspace{1em}

\noindent
\textbf{Date:} October 15, 2026 \\
\textbf{Subject:} Submission of Original Research Manuscript: \textit{Steering Through the Trees: Physical Lie Vector Field Evasion of Tree Ensembles in Network Intrusion Detection}

\vspace{1.5em}

\noindent
Dear Editor-in-Chief and Editorial Board Members,

\vspace{0.8em}

\noindent
We are pleased to submit our original research manuscript entitled \textbf{``Steering Through the Trees: Physical Lie Vector Field Evasion of Tree Ensembles in Network Intrusion Detection''} for consideration for publication in \textit{Computers \& Security}.

\vspace{0.8em}
\noindent
\textbf{Context and Motivation:} \\
Tree ensemble models---most notably Gradient Boosted Decision Trees (XGBoost, LightGBM, CatBoost) and Random Forests---have become the gold standard in production Network Intrusion Detection Systems (NIDS), frequently reporting near-perfect detection metrics ($>99.8\%$ accuracy and ROC-AUC scores exceeding $99.9\%$). However, existing adversarial machine learning (AML) research exhibits a fundamental flaw: unconstrained feature-space perturbations treat tabular flow records as arbitrary continuous coordinates in $\mathbb{R}^d$, producing synthetic records that violate basic network conservation laws, protocol state machines, and transport layer invariants (yielding real-world Physical Feasibility Rates as low as $17.8\%$). Conversely, packet-level heuristic mutators (genetic algorithms, PSO) act as black-box trial-and-error routines requiring thousands of queries without exploiting tree partitioning geometry.

\vspace{0.8em}
\noindent
\textbf{Key Contributions and Novelty:} \\
In this manuscript, we present \textbf{Flow-Field Evasion (FFE)}, the first non-heuristic adversarial evasion framework founded on the differential geometry of network flows:
\begin{enumerate}
    \item \textbf{Lie Cone of Infinitesimal Flow Generators:} We formally define four infinitesimal vector fields corresponding to elementary packet-level transformations: temporal dilation, payload padding, chaff packet injection, and TCP window morphing.
    \item \textbf{Universal Conservation Invariant Preservation:} We prove mathematically via the Carath{\'e}odory existence theorem and algebraic cancellation that integrating trajectories along this positive Lie Cone strictly preserves flow conservation laws and TCP/IP state machines ($\Psi(x) = \mathbf{0}$) by construction.
    \item \textbf{Comprehensive Benchmarking \& Zero-Defect Performance:} Evaluated across three benchmark corpora (CSE-CIC-IDS2018, UNSW-NB15, and CIDDS-001) against CUDA-accelerated XGBoost and multithreaded LightGBM across five attack families, FFE achieves a \textbf{100.0\% Attack Success Rate} with \textbf{100.0\% Physical Feasibility}, requiring an average of only \textbf{6.8--17.8 model queries} and an execution latency under \textbf{70 ms} per flow.
    \item \textbf{Defensive Hardening Roadmap:} We propose two defenses---Malleability-Penalized Tree Splitting (MPS) and Lie-Guided Adversarial Training (LGAT)---which systematically suppress evasion success from $100.0\%$ down to $11.0\%$ under joint deployment without sacrificing clean detection accuracy.
\end{enumerate}

\vspace{0.8em}
\noindent
\textbf{Fit with \textit{Computers \& Security}:} \\
Given \textit{Computers \& Security}'s longstanding leadership in practical network security, intrusion detection systems, and rigorous empirical evaluations of machine learning defenses, we believe our findings will be of direct, significant interest to the journal's readership. 

\vspace{0.8em}
\noindent
\textbf{Declarations:}
\begin{itemize}
    \item This manuscript represents original work that has not been published previously and is not under consideration for publication elsewhere.
    \item All authors have reviewed and approved the manuscript for submission.
    \item The authors confirm that there are no known conflicts of interest associated with this publication.
    \item We intend to publish under the traditional subscription model (\$0 APC).
\end{itemize}

\noindent
Thank you very much for your time and consideration of our work.

\vspace{1.5em}

\noindent
Sincerely,

\vspace{0.5em}
\noindent
\textbf{The Authors} \\
Corresponding Author: \href{mailto:anonymous@institution.org}{anonymous@institution.org}

\end{document}
